\section{Ethics and Privacy}\label{sec:ethics}

\paragraph{Bystanders.} A head-worn camera sees people who did not consent to being recorded, and bystanders' concerns about camera glasses are well documented~\cite{denning2014situ,hoyle2014privacy}. SpatialNet runs on the glasses, so localisation needs no frame to leave the device, and a downstream step can be limited to the target crop and to moments when the wearer points. SpatialNet is trained to ignore other people's hands; it does not detect, identify or track people. On-device processing still does not remove the need for visible recording indicators; the test unit has user- and bystander-facing LEDs.

\paragraph{Data.} EgoPoint is used under the access terms of its authors' release, for research only. We do not redistribute its images or our converted labels; our conversion script regenerates the labels from the original release. The images shown in this paper come from its test and training splits. \tbd{Real-frame check (Section~\ref{sec:realframes}): state who recorded the frames and how they consented, the ethics approval or exemption, how long the frames are kept, and how bystanders were kept out of view or blurred.}

\paragraph{Availability.} We will release the code, the dataset conversion script, the trained weights and the per-image evaluation outputs, from which every table can be recomputed, upon publication.

\section{Conclusion}\label{sec:conclusion}
SpatialNet localises the object a smart-glasses wearer points at from one RGB frame, with one forward pass of a 0.72M-parameter CNN and no prompt, language model or external pose model, and it runs on a pre-production pair of glasses in 60--80\,ms per frame. It finds the wearer's hand and fingertip reliably and the pointed-at object with a P@0.5 of 58.9\%, but only 25.4\% for large-gap targets. That is well below large GPU models, and compact detectors of the same compute, even with the same ImageNet pretraining, draw tighter boxes, although SpatialNet places its target on the labelled object more often, and an off-the-shelf recogniser names the target from its crops as often as from theirs, except when the target lies far from the fingertip. SpatialNet is also smaller and returns the hand, fingertip and direction from the same pass. Its pasted-hand training keeps it stable when synthetic hands are in view, and does the same for a generic detector. Its instance-anchored heads make the hand box and the direction more precise than a plain model of the same size, and input resolution is the largest further lever. An explicit ray prior learns mainly fingertip proximity when most targets overlap the fingertip, and neither a forward constraint nor re-weighting large-gap targets corrects it. Finally, at least 41\% of the EgoPoint test images are COCO photographs; COCO pretraining did not exploit this in our runs, but benchmarks on this dataset should report it. The main open problems are pointing scenes that EgoPoint does not contain, real bystander hands, and large-gap targets along ambiguous rays. We expect progress on them to come from richer egocentric pointing data at least as much as from larger models.
